% ═══════════════════════════════════════════════════════════════════════════════
%  CAPÍTULO — Bridge
% ═══════════════════════════════════════════════════════════════════════════════

\chapter{Bridge}

% ─── Situación Problema ──────────────────────────────────────────────────────
\section{Situación Problema}

El panel de anfitriones de Airbnb permite exportar informes sobre el
rendimiento de sus propiedades. Existen distintos \emph{tipos de informe}
que un anfitrión puede solicitar: resumen de ingresos, estadísticas de
ocupación e historial de reseñas. A su vez, cada informe puede
descargarse en distintos \emph{formatos de salida}: PDF para archivar,
CSV para importar en hojas de cálculo y JSON para integraciones con
herramientas externas de gestión.

Si se modela cada combinación como una clase independiente
---\texttt{InformeIngresosPDF}, \texttt{InformeOcupacionCSV},
\texttt{InformeResenasJSON}, etc.--- la jerarquía crece de forma
combinatoria: tres tipos de informe por tres formatos producen nueve
clases, y agregar un nuevo formato obliga a crear una nueva subclase
por cada tipo de informe existente. La lógica de generación de datos
queda duplicada entre clases que en esencia comparten el mismo proceso
de recopilación, diferenciándose solo en cómo serializan el resultado.

% ─── Solución ────────────────────────────────────────────────────────────────
\section{Solución}

El patrón Bridge separa la jerarquía de \emph{tipos de informe}
---la abstracción--- de la jerarquía de \emph{formatos de exportación}
---la implementación--- conectándolas mediante composición. La clase
\texttt{InformeAnfitrion} contiene una referencia a un objeto
\texttt{ExportadorFormato}; los tipos concretos de informe
(\texttt{InformeIngresos}, \texttt{InformeOcupacion},
\texttt{InformeResenas}) extienden la abstracción y se encargan de
recopilar los datos propios de cada análisis, mientras que los
exportadores concretos (\texttt{ExportadorPDF}, \texttt{ExportadorCSV},
\texttt{ExportadorJSON}) implementan exclusivamente la serialización
del resultado en su formato correspondiente.

Ambas dimensiones evolucionan de forma completamente independiente:
añadir un nuevo tipo de informe no requiere replicar la lógica de
exportación, e incorporar un nuevo formato no impone cambios en ningún
informe existente. El número total de clases crece linealmente en lugar
de de forma combinatoria, y cada clase mantiene una única responsabilidad
bien definida.

% ─── Diagrama de clases ─────────────────────────────────────────────────────
\begin{figure}[H]
    \centering
    % \includegraphics[width=\textwidth]{imagenes/abstract_factory_diagrama.png}
    \caption{Diagrama de clases del patrón Abstract Factory}
\end{figure}


% ─── Roles de las Clases ────────────────────────────────────────────────────
\section{Roles de las Clases}

% Explicar la responsabilidad de cada clase/interfaz

\begin{itemize}
    \item \textbf{AbstractFactory:}
    \item \textbf{ConcreteFactory:}
    \item \textbf{AbstractProduct:}
    \item \textbf{ConcreteProduct:}
    \item \textbf{Client:}
\end{itemize}


% ─── Main ───────────────────────────────────────────────────────────────────
\section{Main}

% Explicar qué realiza el programa principal

\begin{lstlisting}[language=Java, caption=NombreClase]
 
\end{lstlisting}


% ─── Salida del Main ────────────────────────────────────────────────────────
\section{Salida del Main}

\begin{lstlisting}[language=Java, caption=NombreClase]
 
\end{lstlisting}


% ─── Clases ─────────────────────────────────────────────────────────────────
\section{Clases}

% ─── Clase 1 ────────────────────────────────────────────────────────────────
\subsection{NombreClase}

\begin{lstlisting}[language=Java, caption=NombreClase]
 
\end{lstlisting}


% ─── Clase 2 ────────────────────────────────────────────────────────────────
\subsection{NombreClase}

\begin{lstlisting}[language=Java, caption=NombreClase]
 
\end{lstlisting}


% ─── Clase 3 ────────────────────────────────────────────────────────────────
\subsection{NombreClase}

\begin{lstlisting}[language=Java, caption=NombreClase]
 
\end{lstlisting}


% ─── Conclusiones ───────────────────────────────────────────────────────────
\section{Conclusiones}

% Explicar:
% - Ventajas del patrón
% - Cuándo usarlo
% - Beneficios obtenidos
% - Posibles desventajas
% - Relación con principios SOLID